\documentclass[11pt]{article}
\usepackage[margin=1in]{geometry}
\usepackage[T1]{fontenc}
\usepackage{lmodern}
\usepackage{microtype}
\usepackage{amsmath,amssymb,amsthm,mathtools,mathrsfs}
\usepackage{booktabs,array,enumitem}
\usepackage[dvipsnames]{xcolor}
\usepackage[hidelinks]{hyperref}
\usepackage[nameinlink,capitalise]{cleveref}

\newtheorem{theorem}{Theorem}[section]
\newtheorem{proposition}[theorem]{Proposition}
\newtheorem{lemma}[theorem]{Lemma}
\newtheorem{corollary}[theorem]{Corollary}
\newtheorem{criterion}[theorem]{Criterion}
\theoremstyle{definition}
\newtheorem{definition}[theorem]{Definition}
\newtheorem{question}[theorem]{Question}
\theoremstyle{remark}
\newtheorem{remark}[theorem]{Remark}

\newcommand{\F}{\mathbb F}
\newcommand{\N}{\mathbb N}
\newcommand{\Z}{\mathbb Z}
\newcommand{\K}{\mathcal K}
\newcommand{\SSing}{\mathcal S}
\newcommand{\Iness}{\mathcal{IN}}
\newcommand{\B}{\mathcal B}
\newcommand{\Cal}{\operatorname{Cal}}
\newcommand{\ind}{\operatorname{ind}}
\newcommand{\ran}{\operatorname{ran}}
\newcommand{\codim}{\operatorname{codim}}
\newcommand{\Id}{I}
\newcommand{\eps}{\varepsilon}
\newcommand{\sq}{\oplus}

\title{Square Stability versus Hyperplane Stability in Banach Spaces\\
\large Audit of the supplied Cuntz-pair obstruction, literature status, and new reductions}
\author{Research report prepared from the supplied note and source paper}
\date{22 June 2026}

\begin{document}
\maketitle

\begin{abstract}
Albiac and Ansorena ask whether there exists a Banach space $X$ such that
$X\cong X\oplus X$ but $X\not\cong X\oplus\F$.  The supplied note does not
solve this problem; it proves a no-go statement for one attempted
Gowers--Maurey spread construction.  I verify that its conclusion is correct,
identify a small definitional gap in the displayed proof, and give a proof that
works even under the weaker definition written in the note.  I also extend the
obstruction from binary to every fixed finite arity and prove an address-rigidity
result for doubly commuting Cuntz spread families.

A literature search through 22 June 2026 found no published or preprint proof or
counterexample.  In particular, the March 2026 source paper still states the
problem as open, and a February 2026 paper still calls the hyperplane problem
for the square-stable Kalton--Peck space $Z_2$ open.  The report develops several
additional reductions.  The sharpest are: a $K_0$ characterization in terms of
the rank-one idempotent; a reduction to the old Gowers--Maurey spaces $S_m$;
a conditional even-index theorem for $Z_2$; and a realification theorem saying
that a square-stable complex space essentially incomparable with its conjugate
would immediately yield a real counterexample.  These observations materially
narrow the problem, but they do not close the final gap.  New deductions in this
report have not been peer reviewed.
\end{abstract}

\tableofcontents

\section{Outcome}

Let $X$ be an infinite-dimensional Banach space over $\F=\mathbb R$ or
$\mathbb C$.  The question is
\begin{equation}\label{eq:main-question}
 X\cong X^2:=X\oplus X
 \qquad\text{and}\qquad
 X\not\cong X\oplus\F.
\end{equation}
The first property will be called \emph{square stability}, and the negation of
the second is \emph{hyperplane nonstability}.

The conclusions of this report are as follows.
\begin{enumerate}[label=\textup{(\arabic*)},leftmargin=2.2em]
\item The supplied finite-overlap obstruction is mathematically correct.  As
written, however, its proof uses closure of the spread family under products
and adjoints, a clause present in the standard Gowers--Maurey definition but
omitted from the note's displayed definition.  \Cref{thm:binary-repair} repairs
the proof without adding that clause.
\item The obstruction extends to exact $d$-ary Cuntz spread families for every
fixed $d\ge2$; see \Cref{thm:dary}.  It rules out an infinite reserve of fresh
exact decompositions inside one proper family, but it does \emph{not} settle the
original single dyadic-tree strategy proposed by Gowers and Maurey.
\item No full answer was found in the literature.  The problem was explicitly
open in March 2026 \cite{AlbiacAnsorena2026}; the closely related $Z_2$
hyperplane problem was explicitly open in February 2026
\cite{CarreraFerenczi2026}.
\item Several rigorous routes reduce the problem to concrete operator-algebraic
statements.  The most promising currently appear to be the dyadic
Gowers--Maurey index-transfer problem and triangularization of Fredholm
operators on $Z_2$.
\end{enumerate}

\section{Fredholm and $K$-theoretic reformulations}

\subsection{The index subgroup}

For a Fredholm operator $T\in\B(X)$ use the convention
\[
 \ind T=\dim\ker T-\codim(TX;X).
\]
The set of all such indices is an additive subgroup of $\Z$.

\begin{proposition}\label{prop:index-subgroup}
There is a unique $h(X)\in\N_0$ such that
\[
 \{\ind T:T\in\B(X)\text{ is Fredholm}\}=h(X)\Z,
\]
where $0\Z=\{0\}$.  Moreover, for $n\ge1$,
\[
 X\cong X\oplus\F^n
 \quad\Longleftrightarrow\quad
 n\in h(X)\N.
\]
In particular, $X$ is hyperplane stable if and only if a Fredholm operator on
$X$ has index $1$ (equivalently, index $-1$).
\end{proposition}

\begin{proof}
Fredholm indices add under composition, change sign on a parametrix, and include
$0$, so they form a subgroup of $\Z$.  Every subgroup is $h\Z$.  If $T$ has
index $n\ge0$, choose complements to its kernel and range.  The resulting
finite-dimensional bookkeeping gives $X\cong X\oplus\F^n$.  Conversely, if
$J:X\to X\oplus\F^n$ is an isomorphism and $P:X\oplus\F^n\to X$ is the
coordinate projection, then $PJ$ is Fredholm of index $n$.
\end{proof}

Thus the problem asks for a square-stable space whose index subgroup is either
$\{0\}$ or $d\Z$ for some $d\ge2$.

\subsection{A sharp $K_0$ test}

This subsection is over the complex field; an analogous statement can be made
with real $K$-theory.  Let $q\in\B(X)$ be any rank-one idempotent.  Write
$[q]_0,[I]_0\in K_0(\B(X))$ for their classes.

\begin{theorem}[rank-one criterion]\label{thm:k0-criterion}
Assume that $X\cong X^2$.  Then
\[
 X\cong X\oplus\mathbb C
 \quad\Longleftrightarrow\quad
 [q]_0=0\quad\text{in }K_0(\B(X)).
\]
Consequently, every complex solution of \eqref{eq:main-question} must satisfy
\[
 [I]_0=0\quad\text{but}\quad[q]_0\ne0.
\]
In particular, $K_0(\B(X))$ must be nontrivial even though its unit class
vanishes.
\end{theorem}

\begin{proof}
If $X\cong X\oplus\mathbb C$, then $I\oplus q$ and $I$ have isomorphic ranges,
so $[q]_0=0$.

Conversely, $[q]_0=0$ means that, after stabilization, there is an
idempotent $e\in M_n(\B(X))$ for which $q\oplus e$ is equivalent to $e$.
Equivalence of idempotents in matrix algebras over $\B(X)$ is equivalence of
their ranges.  If $Y=\operatorname{ran}e$, this gives
\[
 \mathbb C\oplus Y\cong Y.
\]
The space $Y$ is complemented in $X^n$.  Square stability implies
$X^n\cong X$, so, after transporting the projection, write $X\cong Y\oplus Z$.
It follows that
\[
 X\oplus\mathbb C\cong (Y\oplus\mathbb C)\oplus Z
 \cong Y\oplus Z\cong X.
\]

Finally, $X\cong X^2$ gives $[I]_0=2[I]_0$, hence $[I]_0=0$.
\end{proof}

Under square stability, the same argument gives
\[
 n[q]_0=0\quad\Longleftrightarrow\quad X\cong X\oplus\mathbb C^n.
\]
Thus $[q]_0$ has order $h(X)$ when $h(X)>0$, and infinite order when
$h(X)=0$.  This agrees with Laustsen's Fredholm--$K$-theoretic framework
\cite{Laustsen1999}.

\begin{corollary}\label{cor:calkin-k1}
Let $X$ be a complex Banach space with $X\cong X^2$.  If
\[
 K_1(\Cal(X))=0,
 \qquad \Cal(X):=\B(X)/\K(X),
\]
then $X$ is not isomorphic to a hyperplane.
\end{corollary}

\begin{proof}
The Fredholm index factors through the stable homotopy class of the invertible
coset in $\Cal(X)$, hence through $K_1(\Cal(X))$.  If this group is zero, every
Fredholm operator has index zero.  Apply \Cref{prop:index-subgroup}.
\end{proof}

This gives a clean construction blueprint: realize square stability by an exact
binary Cuntz family while forcing a Calkin algebra with trivial $K_1$.  The
known constructions achieve these two requirements separately, but no known
construction controls both simultaneously.

\subsection{Structural restrictions on a counterexample}

\begin{proposition}[primary obstruction]\label{prop:primary}
If $X$ is primary and $X\cong X^2$, then $X\cong X\oplus\F$.
\end{proposition}

\begin{proof}
The space $X\oplus\F$ is isomorphic to a complemented subspace of $X^2$, hence
of $X$.  Write
\[
 X\cong (X\oplus\F)\oplus Z.
\]
By primarity, either $X\oplus\F\cong X$, and there is nothing to prove, or
$Z\cong X$.  In the latter case
\[
 X\cong(X\oplus\F)\oplus X\cong X^2\oplus\F\cong X\oplus\F.
\]
\end{proof}

Kalton proved stronger results for countably primary spaces with unconditional
bases \cite{Kalton1999}.  Thus a solution must be decidedly nonprimary.

\begin{proposition}[shape of a hypothetical hyperplane]\label{prop:hyperplane-shape}
Assume $X\cong X^2$ and let $H$ be a hyperplane of $X$.  Then
\[
 H\cong X\oplus H.
\]
Hence $X$ and $H$ complementably embed into one another.  If $X\not\cong H$,
they form a particularly rigid failure of the Banach-space
Schroeder--Bernstein principle.
\end{proposition}

\begin{proof}
Pull $H$ back through an isomorphism $X\oplus X\to X$.  The inverse image is
the kernel of a nonzero functional $(f,g)$ on $X\oplus X$.  If, say, $g\ne0$,
splitting $X=\ker g\oplus\F$ identifies this kernel with $X\oplus\ker g$.
All hyperplanes of $X$ are mutually isomorphic, so $\ker g\cong H$.
\end{proof}

\begin{lemma}[stable-reservoir obstruction]\label{lem:reservoir}
If $X=Y\oplus Z$ and $Y\cong Y\oplus\F$, then $X\cong X\oplus\F$.
\end{lemma}

\begin{proof}
$X\oplus\F\cong(Y\oplus\F)\oplus Z\cong Y\oplus Z=X$.
\end{proof}

Consequently, a counterexample contains no complemented hyperplane-stable
space.  This eliminates the standard countable $\ell_p$- and $c_0$-sum
stabilizations: they contain a complemented scalar $\ell_p$ or $c_0$ diagonal.
It also explains why tensoring a hyperplane-nonstable space with a classical
stable factor is counterproductive whenever that factor survives complementedly.

\section{Audit and repair of the supplied spread argument}

\subsection{The omitted clause}

For infinite $A=\{a_1<a_2<\cdots\}$ and
$B=\{b_1<b_2<\cdots\}$, the Gowers--Maurey spread $S_{A,B}$ is defined by
$S_{A,B}e_{a_k}=e_{b_k}$ and $S_{A,B}e_n=0$ for $n\notin A$.

In the standard definition \cite[Section 3]{GowersMaurey1997}, a proper family
of spreads is
\begin{enumerate}[label=\textup{(\roman*)}]
\item closed under composition and formal adjoints; and
\item subject to the finite-overlap condition: for two distinct matrix
positions, only finitely many family members are nonzero at both positions.
\end{enumerate}
The supplied note displays only (ii).  Its proof then asserts that
$P_A=R R^*$ belongs to the family.  That is automatic from (i), but not from the
weaker displayed definition.  Thus the proof is valid under the actual standard
definition, while under the note's literal definition it has a small gap.

The following proof removes the gap entirely.

\subsection{A repaired binary theorem}

Call a family \emph{weakly proper} if it satisfies only the finite-overlap
condition.  An exact binary Cuntz spread pair is a pair $(R_0,R_1)$ of spreads
with
\[
 R_i^*R_j=\delta_{ij}I,
 \qquad R_0R_0^*+R_1R_1^*=I.
\]
Each $R_i$ is therefore the unique spread from all of $\N$ onto an infinite set
$A_i$, where $\N=A_0\sqcup A_1$.

\begin{theorem}[repaired supplied result]\label{thm:binary-repair}
A weakly proper family of spreads contains only finitely many exact binary
Cuntz pairs.
\end{theorem}

\begin{proof}
Suppose there are infinitely many pairs.  Since a fixed unordered partition
$\N=A_0\sqcup A_1$ gives at most two labelled pairs, there are infinitely many
distinct unordered binary partitions.

Choose three coordinates $a<b<c$.  In each partition, two of these three
coordinates lie in one cell.  Passing to an infinite subfamily, a fixed pair,
say $a,b$, lies in a cell $B_m$.  For binary partitions the cells $B_m$ may be
assumed distinct: one cell determines its complement and hence the unordered
partition.

Let $R_m$ be the member of the corresponding Cuntz pair with range $B_m$.
Define
\[
 r_m(t)=|B_m\cap\{1,\dots,t\}|.
\]
Because $R_m$ is the increasing spread from $\N$ onto $B_m$,
\[
 R_m e_{r_m(a)}=e_a,
 \qquad
 R_m e_{r_m(b)}=e_b.
\]
There are only $a$ possible values of $r_m(a)$ and only $b$ possible values of
$r_m(b)$.  A second pigeonhole argument therefore gives an infinite subfamily
on which both ranks are fixed.  Infinitely many distinct spreads are then
nonzero at the same two distinct matrix positions
\[
 (a,r(a))\quad\text{and}\quad(b,r(b)),
\]
contradicting weak properness.
\end{proof}

This proves that the conclusion of the supplied packet is right even with its
weaker definition.  The original projection proof is shorter if standard
closure is restored.

\subsection{Finite arity}

The binary argument has a useful finite-arity extension.

\begin{lemma}[partition lemma]\label{lem:partition}
Fix $d\ge2$.  Let $\mathscr P$ be an infinite family of distinct partitions of
$\N$ into $d$ infinite cells.  Then there exist distinct coordinates $a,b$ and
infinitely many distinct cells, drawn from partitions in $\mathscr P$, each
containing both $a$ and $b$.
\end{lemma}

\begin{proof}
Induct on $d$.  The case $d=2$ is the three-coordinate argument used above.
For general $d$, choose $d+1$ coordinates.  In every partition two lie in a
common cell, so a fixed pair lies in a cell $C_P$ for infinitely many
partitions.  If infinitely many $C_P$ are distinct, we are done.  Otherwise,
pass to an infinite subfamily with one fixed cell $C$.  Removing $C$ leaves
infinitely many distinct partitions of $\N\setminus C$ into $d-1$ infinite
cells.  Apply the induction hypothesis.
\end{proof}

\begin{theorem}[finite-arity no-go]\label{thm:dary}
For every fixed $d\ge2$, a weakly proper family of spreads contains only
finitely many exact $d$-ary Cuntz spread families
$(R_1,\dots,R_d)$ satisfying
\[
 R_i^*R_j=\delta_{ij}I,
 \qquad \sum_{i=1}^dR_iR_i^*=I.
\]
\end{theorem}

\begin{proof}
Infinitely many such families give infinitely many distinct unordered
$d$-cell partitions, up to the finite number $d!$ of labelings.  Apply
\Cref{lem:partition} to obtain a fixed pair $a,b$ in infinitely many distinct
range cells $B_m$.  Select the corresponding spread $R_m:\N\to B_m$ and fix
the two finite-valued ranks $|B_m\cap[1,a]|$ and $|B_m\cap[1,b]|$ on an
infinite subfamily.  This violates finite overlap exactly as in
\Cref{thm:binary-repair}.
\end{proof}

\subsection{Address rigidity and double commutation}

Under the full standard definition, properness gives an additional obstruction.

\begin{proposition}[address rigidity]\label{prop:address}
Let $\mathscr S$ be a standard proper family and let $(R_0,R_1)\subset\mathscr S$
be an exact binary Cuntz pair.  For a word $w$ in $\{0,1\}$ put
$R_w=R_{w_1}\cdots R_{w_n}$ and $P_w=R_wR_w^*$.  The binary address map
\[
 a:\N\longrightarrow\{0,1\}^{\N},
 \qquad
 a(x)|_n=w\ \Longleftrightarrow\ P_we_x=e_x,
\]
is injective.
\end{proposition}

\begin{proof}
At every level the projections $(P_w)_{|w|=n}$ are pairwise disjoint and sum to
$I$, so every coordinate has a unique address.  Along a branch the cylinder
projections are distinct, since each cylinder splits into two nonzero child
cylinders.  If distinct coordinates $x,y$ had the same address, then the
infinitely many projections $P_{a(x)|_n}\in\mathscr S$ would all be nonzero at
the two diagonal positions $(x,x)$ and $(y,y)$, contrary to properness.
\end{proof}

\begin{corollary}[no second doubly commuting pair]\label{cor:double-commuting}
If a total spread $T\in\mathscr S$ commutes with every $R_i$ and $R_i^*$, then
$T=I$.  Consequently, $\mathscr S$ cannot contain a second exact Cuntz pair
that doubly commutes with the first one across generators.
\end{corollary}

\begin{proof}
The commutation assumptions imply that $T$ commutes with every $P_w$.  If
$Te_x=e_y$, then $x$ and $y$ belong to exactly the same cylinder ranges, so
\Cref{prop:address} gives $x=y$.  Hence $T=I$.  Two members of a second exact
pair cannot both equal $I$.
\end{proof}

Ordinary commutation without the adjoint relations is weaker and is not covered
by this corollary.

\section{What the literature currently gives}

\subsection{Status as of 22 June 2026}

The source paper states the question explicitly and describes it as open
\cite[Question 2.15]{AlbiacAnsorena2026}.  Searches were made for the exact
question, for the phrases ``isomorphic to its square'' and ``hyperplane'', for
updates on the Gowers--Maurey $S_m$ and $H_m$ spaces, and for the hyperplane
problem for $Z_2$.  No later proof or counterexample was found.

Two independent recent checkpoints reinforce that conclusion.
\begin{itemize}[leftmargin=1.6em]
\item Castillo, Gonzalez and Pino state that $Z_2$ is isomorphic to its
codimension-two subspaces while its codimension-one problem remains unknown
\cite{CastilloGonzalezPino2025}.
\item Carrera and Ferenczi again call the $Z_2$ hyperplane problem open in a
February 2026 preprint \cite{CarreraFerenczi2026}.
\end{itemize}
Since $Z_2\cong Z_2^2$, a negative answer to its hyperplane problem would solve
\eqref{eq:main-question} immediately.

\subsection{The two original Gowers--Maurey families}

Gowers and Maurey constructed two relevant families \cite{GowersMaurey1997}.
For each $m\ge2$ there is a space $H_m$ whose finite-codimensional isomorphism
classes are periodic modulo $m$; in particular $h(H_m)=m$.  Thus:
\[
 \boxed{\text{If any }H_m\text{ is square stable, the problem is solved.}}
\]
The square stability of these spaces is still unknown in the source paper.

They also constructed spaces $S_m$ with periodic powers:
\[
 S_m^a\cong S_m^b\quad\Longleftrightarrow\quad a\equiv b\pmod m.
\]
Therefore $Y_m:=S_m^m$ is automatically square stable, since
$S_m^{2m}\cong S_m^m$.

\begin{proposition}[reduction to $S_m$]\label{prop:Sm-reduction}
If $S_m$ is not hyperplane stable for some $m\ge2$, then $S_m^m$ solves
\eqref{eq:main-question}.
\end{proposition}

\begin{proof}
Only hyperplane nonstability needs proof.  Suppose $S_m^m$ were hyperplane
stable.  Adding one copy of $S_m$ gives
\[
 S_m^{m+1}\oplus\F\cong S_m^{m+1}.
\]
But $S_m^{m+1}\cong S_m$ by periodicity, so $S_m\cong S_m\oplus\F$, a
contradiction.
\end{proof}

Gowers and Maurey conjectured the stronger assertion $h(S_m)=0$.  Thus their
old conjecture would answer the 2026 question at once.

\section{The single-tree Gowers--Maurey route}

\subsection{What the 1997 argument proves}

In their ternary construction, Gowers and Maurey use right-append operators
$S_0,S_1,S_2$ and left-prefix operators $T_0,T_1,T_2$ on the rooted ternary
tree.  They commute, satisfy
\[
 S_i^*S_j=T_i^*T_j=\delta_{ij}I,
 \qquad
 \sum_iS_iS_i^*=\sum_iT_iT_i^*=I-P,
\]
where $P$ is the rank-one root projection.  Modifying the right shifts at the
root produces an exact ternary Cuntz family and hence a space $X\cong X^3$.

On an auxiliary $\ell_1$ tree space they define
\[
 \Phi(U)=\sum_{i=0}^2 T_iUT_i^*.
\]
For operators in the tree algebra, $\Phi(U)-U$ has finite rank, while
$\Phi(U)$ is equivalent to three copies of $U$ plus a root defect.  This gives
for a rectangular Fredholm matrix $U:Y^q\to Y^p$ the identity
\[
 2\ind U=p-q.
\]
In particular, every scalar Fredholm operator in the auxiliary algebra has
index zero, and no rectangular Fredholm operator exists when $p-q$ is odd.
This is the engine behind $X\cong X^3$ but $X\not\cong X^2$.

The authors then write that it is likely every Fredholm operator on the actual
space has index zero and that a dyadic tree may yield a space isomorphic to its
square but not to its hyperplanes \cite[Section 4.4]{GowersMaurey1997}.
That sentence is essentially the present problem in embryonic form.

\subsection{Why the supplied no-go does not stop this route}

The supplied result excludes a countable reserve of fresh exact Cuntz spread
pairs in one proper family.  The 1997 proposal does not require such a reserve:
it asks whether one dyadic family, together with the auxiliary tree algebra,
controls all Fredholm indices.  Thus the supplied theorem is a valid route
obstruction, not a negative answer to the dyadic program.

\subsection{The precise missing transfer}

The Gowers--Maurey theorem yields a homomorphism from $\B(X)$ into a completion
of the spread algebra modulo a seminorm-zero ideal.  A further homomorphism
lands in an auxiliary tree quotient.  This is strong enough to detect the
rectangular obstruction used to distinguish powers.  It is not known to retain
the finite-dimensional kernel/cokernel defect of a Fredholm endomorphism on
$X$.  In particular, finite-rank operators disappear in the quotient, exactly
where the integer index is stored.

A complete dyadic proof would follow from any one of the following types of
statement.
\begin{enumerate}[label=\textup{(\alph*)},leftmargin=2em]
\item a natural morphism of extensions that sends a rank-one operator on $X$ to
the root rank-one operator in the auxiliary tree algebra;
\item injectivity, on the relevant $K_1$ classes, of the symbol map from the
Calkin algebra of $X$ to the auxiliary quotient; or
\item a direct amplification on $X$ that is simultaneously a finite-rank (or
inessential) perturbation of the identity on the whole approximating algebra
and doubles Fredholm index.
\end{enumerate}
The finite-overlap theorem explains why option (c) cannot be implemented by an
infinite list of fresh exact spread pairs.  Options (a) and (b) remain viable.

\section{The Kalton--Peck candidate $Z_2$}

\subsection{Known operator structure}

The Kalton--Peck space fits into a singular exact sequence
\[
 0\longrightarrow\ell_2\xrightarrow{i}Z_2\xrightarrow{p}\ell_2
 \longrightarrow0.
\]
It is isomorphic to its square and to its codimension-two subspaces.  Recent
work proves, among other things, that on $Z_2$:
\begin{itemize}[leftmargin=1.6em]
\item strictly singular, strictly cosingular and inessential operators coincide;
\item the product of two strictly singular operators is compact; and
\item every semi-Fredholm operator has complemented kernel and range.
\end{itemize}
See \cite{CastilloGonzalezPino2025}.

An operator can be represented schematically by
\[
 T=\begin{pmatrix}\alpha&\beta\\ \delta&\gamma\end{pmatrix}.
\]
If $\delta=0$, the operator is upper triangular relative to the exact sequence;
then $\alpha,\gamma\in\B(\ell_2)$ and
\[
 \alpha-\gamma\in\K(\ell_2).
\]
This almost forces parity of the index.

\subsection{A conditional parity theorem}

\begin{theorem}[upper-triangular parity]\label{thm:z2-parity}
Let $T\in\B(Z_2)$ be Fredholm.  Suppose that $T$ is an inessential (in
particular, strictly singular) perturbation of an upper triangular operator
\[
 U=\begin{pmatrix}\alpha&\beta\\0&\gamma\end{pmatrix}.
\]
Then
\[
 \ind T=2\ind\alpha,
\]
and in particular $\ind T$ is even.
\end{theorem}

\begin{proof}
Inessential perturbations preserve Fredholmness and index, so it suffices to
consider $U$.  The upper triangular form gives a commutative morphism of the
short exact sequence defining $Z_2$, with $\alpha$ on the kernel and $\gamma$
on the quotient.  The algebraic snake sequence is
\[
0\to\ker\alpha\to\ker U\to\ker\gamma\to
\operatorname{coker}\alpha\to\operatorname{coker}U\to
\operatorname{coker}\gamma\to0.
\]
Since $U$ is Fredholm, $\ker\alpha$ is finite-dimensional and
$\operatorname{coker}\gamma$ is finite-dimensional.  Hence $\gamma$ is lower
semi-Fredholm.  The compact relation $\alpha-\gamma\in\K(\ell_2)$ implies that
$\alpha$ is also lower semi-Fredholm; together with finite-dimensional kernel,
$\alpha$ is Fredholm.  Compact perturbation then makes $\gamma$ Fredholm with
$\ind\gamma=\ind\alpha$.  Taking the alternating sum of dimensions in the
snake sequence gives
\[
 \ind U=\ind\alpha+\ind\gamma=2\ind\alpha.
\]
\end{proof}

\begin{corollary}\label{cor:z2-solve}
If every Fredholm operator on $Z_2$ is an inessential perturbation of an upper
triangular operator, then $Z_2$ is square stable and has index subgroup $2\Z$;
therefore it solves \eqref{eq:main-question}.
\end{corollary}

\begin{proof}
The theorem excludes odd indices, while the known codimension-two isomorphism
supplies an operator of index $2$.
\end{proof}

\subsection{The exact remaining corner}

For every $T$, the lower-left entry can be written $\delta=pTi:\ell_2\to\ell_2$.
The quotient map $p$ is strictly singular, so $\delta$ is strictly singular on
$\ell_2$, hence compact.  Compactness alone is not enough for the known
triangularization theorem.  Castillo--Gonzalez--Pino prove triangularization
when $\delta$ has the stronger smoothing extension
\[
 \delta:\ell_f^*\longrightarrow\ell_f,
\]
in their auxiliary notation.  Proving this extra regularity for Fredholm $T$,
or finding a weaker argument that extracts index parity directly from compact
$\delta$, would settle the problem through $Z_2$.

\section{A complex-conjugate route}

A different route uses the parity forced by realification.  Recall that complex
spaces $E$ and $F$ are \emph{essentially incomparable} when every operator in
both cross directions is inessential.

\begin{theorem}[realification parity]\label{thm:realification}
Let $E$ be a complex Banach space such that
\begin{enumerate}[label=\textup{(\roman*)}]
\item $E\cong E^2$ complex-linearly; and
\item $E$ and its complex conjugate $\overline E$ are essentially incomparable.
\end{enumerate}
Then the underlying real space $E_{\mathbb R}$ is square stable and every real
Fredholm operator on it has even index.  Hence $E_{\mathbb R}$ is not
isomorphic to a real hyperplane.
\end{theorem}

\begin{proof}
Square stability survives restriction of scalars.  Let $T$ be a real Fredholm
operator on $E_{\mathbb R}$.  The complexification of $E_{\mathbb R}$ is
complex-isomorphic to $E\oplus\overline E$.  Relative to this decomposition,
$T_{\mathbb C}$ has the form
\[
 T_{\mathbb C}=
 \begin{pmatrix}A&C\\B&\overline A\end{pmatrix},
\]
where $B:E\to\overline E$ and $C:\overline E\to E$.  Essential
incomparability makes the off-diagonal matrix inessential.  Therefore
$T_{\mathbb C}$ has the same index as
$\operatorname{diag}(A,\overline A)$.  The latter is Fredholm, so $A$ and
$\overline A$ are Fredholm and have equal complex index.  Finally,
\[
 \ind_{\mathbb R}T
 =\ind_{\mathbb C}T_{\mathbb C}
 =\ind_{\mathbb C}A+\ind_{\mathbb C}\overline A
 =2\ind_{\mathbb C}A.
\]
The equality between the real index and the complexified index follows because
the complexification of a finite-dimensional real kernel or cokernel has the
same dimension over $\mathbb C$.
\end{proof}

\begin{corollary}[power-periodic version]\label{cor:complex-power}
Suppose $E^{m+1}\cong E$ for some $m\ge1$ and $E$ is essentially incomparable
with $\overline E$.  Then the real space $(E^m)_{\mathbb R}$ solves
\eqref{eq:main-question}.
\end{corollary}

\begin{proof}
Finite powers preserve essential incomparability.  When $m=1$, the asserted
square stability is exactly $E^2\cong E$.  When $m\ge2$,
\[
 E^{2m}=E^{m-1}\oplus E^{m+1}\cong E^{m-1}\oplus E=E^m.
\]
Apply \Cref{thm:realification} to $E^m$.
\end{proof}

This theorem identifies two complementary pieces already present in the
literature but not yet in one space.
\begin{itemize}[leftmargin=1.6em]
\item Kalton's explicit spaces $Z_\alpha$ are not isomorphic to their complex
conjugates for $\alpha\ne0$ \cite{Kalton1995}, and Wenzel observed that they are
isomorphic to their squares \cite{Wenzel1997}.  What is missing is essential
incomparability; mere nonisomorphism is not enough.
\item Ferenczi constructed complex structures whose conjugates are totally (and
hence essentially) incomparable \cite{Ferenczi2007}, but the underlying spaces
are hereditarily indecomposable and therefore not square stable.
\end{itemize}
A construction combining the square periodicity of the first family with the
cross-operator rigidity of the second would give a full answer by
\Cref{thm:realification}.

The theory of even real Banach spaces gives a related formulation.  An even
space admits a complex structure while its hyperplanes do not
\cite{FerencziGalego2007}.  Thus any square-stable even space would solve the
problem.  Existing even examples are not known to have the required square
stability; the real $Z_2$ is a prominent candidate
\cite{CarreraFerenczi2026}.

\section{Why several natural constructions fail}

\subsection{Finite-dimensional arithmetic is not cancellative}

Let $[X]$ denote the isomorphism class of $X$ in the commutative monoid of
Banach spaces under direct sum.  The desired equations are
\[
 [X]=2[X],
 \qquad [X]\ne[X]+1.
\]
There is no formal cancellation principle in this monoid.  The equality
$[X]=2[X]$ by itself does not absorb the finite-dimensional class $1$; this is
exactly why elementary Eilenberg-swindle arguments stall.

\subsection{Countable sums and tensor stabilizations}

An $\ell_p$-sum $(\bigoplus_nY_n)_p$ with nonzero coordinates contains a
complemented scalar $\ell_p$ diagonal by choosing one vector and one functional
in every coordinate.  The same holds for a $c_0$-sum.  By
\Cref{lem:reservoir}, such a space is hyperplane stable.  Therefore taking
countably many copies of a known hyperplane-nonstable space cannot solve the
problem using a standard unconditional sum.

Likewise, if a tensor or vector-valued construction contains a complemented
copy of a hyperplane-stable factor, the same lemma destroys the desired
obstruction.  A successful stabilization must be highly twisted: it must create
$X\cong X^2$ without leaving any complemented stable reservoir.

\subsection{An exact Cuntz pair is not itself an index-one operator}

An exact pair gives an isomorphism $X^2\cong X$, but a single branch is an
infinite-codimensional range in the natural tree models.  It does not supply a
Fredholm shift with one-dimensional cokernel.  To deduce hyperplane stability
one would need a complemented one-vector-per-level branch or another finite
defect mechanism.  The Gowers--Maurey constructions are designed precisely to
obstruct such complemented selections.

\section{A focused research agenda}

The preceding analysis leaves four concrete targets.

\begin{enumerate}[label=\textbf{\arabic*.},leftmargin=2.2em]
\item \textbf{Dyadic index transfer.}  Construct a morphism of extensions, or a
$K_1$-injective symbol map, connecting the Fredholm index on the dyadic
Gowers--Maurey space to the root-defect index in the auxiliary tree algebra.
This would likely prove all Fredholm indices are zero.

\item \textbf{$Z_2$ regularity.}  Prove that the compact corner
$\delta=pTi$ of every Fredholm operator has enough smoothing to be removed by an
inessential perturbation.  \Cref{thm:z2-parity} would then give index subgroup
$2\Z$.

\item \textbf{Complex-conjugate fusion.}  Produce a square- or power-periodic
complex space essentially incomparable with its conjugate.  The realification
in \Cref{thm:realification,cor:complex-power} would solve the problem without
further Banach-space decomposition arguments.

\item \textbf{$K$-theoretic realization.}  Construct a square-stable $X$ with
$[I]_0=0$ and nonzero rank-one class $[q]_0$, or equivalently with trivial unit
class but nontrivial finite-rank index obstruction.  A realization with
$K_1(\Cal(X))=0$ and an exact lifted binary Cuntz family would be sufficient by
\Cref{cor:calkin-k1}.
\end{enumerate}

\section{Conclusion}

The supplied note contains a valid and useful no-go theorem, after either
restoring the standard closure clause in the definition of a proper spread
family or using the repaired rank-pigeonhole proof above.  Its method extends
to every finite arity and yields an additional double-commutation rigidity
phenomenon.  It does not answer Albiac--Ansorena Question 2.15, and no answer was
found in the literature through 22 June 2026.

The problem can now be phrased sharply in several equivalent or sufficient
forms.  In the complex case, a solution must have a vanishing unit class but a
nonvanishing rank-one class in $K_0(\B(X))$.  The old Gowers--Maurey spaces
reduce the problem to their unresolved hyperplane conjecture.  The Kalton--Peck
space reduces it to Fredholm triangularization and an even-index theorem.  The
complex route reduces it to essential incomparability with the conjugate.  Each
of these is a concrete, falsifiable target; none is presently closed.

\begin{thebibliography}{99}

\bibitem{AlbiacAnsorena2026}
F.~Albiac and J.~L. Ansorena,
\newblock \emph{Unconditional structure of Banach spaces with few operators},
\newblock arXiv:2603.07324 (2026).

\bibitem{CarreraFerenczi2026}
W.~Cuellar Carrera and V.~Ferenczi,
\newblock \emph{On complex structures and uniqueness of algebra norms in Banach spaces},
\newblock arXiv:2602.05045 (2026).

\bibitem{CastilloGonzalezPino2025}
J.~M.~F. Castillo, M.~Gonzalez and R.~Pino,
\newblock Operators on the Kalton--Peck space $Z_2$,
\newblock \emph{Israel J. Math.} \textbf{269} (2025), 185--212.

\bibitem{Ferenczi2007}
V.~Ferenczi,
\newblock Uniqueness of complex structure and real hereditarily indecomposable Banach spaces,
\newblock \emph{Adv. Math.} \textbf{213} (2007), 462--488.

\bibitem{FerencziGalego2007}
V.~Ferenczi and E.~Medina Galego,
\newblock Even infinite-dimensional real Banach spaces,
\newblock \emph{J. Funct. Anal.} \textbf{253} (2007), 534--549.

\bibitem{Gowers1994}
W.~T. Gowers,
\newblock A solution to Banach's hyperplane problem,
\newblock \emph{Bull. London Math. Soc.} \textbf{26} (1994), 523--530.

\bibitem{GowersMaurey1997}
W.~T. Gowers and B.~Maurey,
\newblock Banach spaces with small spaces of operators,
\newblock \emph{Math. Ann.} \textbf{307} (1997), 543--568.

\bibitem{Kalton1995}
N.~J. Kalton,
\newblock An elementary example of a Banach space not isomorphic to its complex conjugate,
\newblock \emph{Canad. Math. Bull.} \textbf{38} (1995), 218--222.

\bibitem{Kalton1999}
N.~J. Kalton,
\newblock A remark on Banach spaces isomorphic to their squares,
\newblock in \emph{Function Spaces (Edwardsville, IL, 1998)}, Contemp. Math. 232,
Amer. Math. Soc., 1999, 211--217.

\bibitem{Laustsen1999}
N.~J. Laustsen,
\newblock $K$-theory for algebras of operators on Banach spaces,
\newblock \emph{J. London Math. Soc.} (2) \textbf{59} (1999), 715--728.

\bibitem{Wenzel1997}
J.~Wenzel,
\newblock A supplement to my paper on real and complex operator ideals,
\newblock \emph{Quaest. Math.} \textbf{20} (1997), 663--665.

\end{thebibliography}

\end{document}
